% !TeX program = xelatex
% !TeX encoding = UTF-8
% 《模空间的几何与量子化》中文译本（试点：前言 + 第 2 章）
% 底本：Geometry and Quantization of Moduli Spaces,
%       Advanced Courses in Mathematics - CRM Barcelona, Birkhäuser 2016.
% 编译：xelatex main.tex （三遍，见 工程文件/build.sh）
\documentclass[UTF8, a4paper, 11pt, openany]{ctexbook}

% ---------- 宏包 ----------
\usepackage{amsmath, amssymb, amsthm, mathtools}
\usepackage{fontspec}
\usepackage{geometry}
\usepackage{fancyhdr}
\usepackage{enumitem}
\usepackage{xcolor}
\usepackage{tikz}
\usepackage{tikz-cd}
\usepackage[colorlinks=true, linkcolor=blue!60!black, citecolor=green!45!black,
            urlcolor=blue!60!black, bookmarksnumbered=true,
            pdftitle={模空间的几何与量子化（中文译本）},
            pdfauthor={Vladimir Fock, Andrey Marshakov, Florent Schaffhauser, Constantin Teleman, Richard Wentworth}]{hyperref}

% ---------- 字体 ----------
\setCJKmainfont{Noto Serif SC}[AutoFakeBold=1.5, AutoFakeSlant=0.2]
\setCJKsansfont{Noto Sans SC}[AutoFakeBold=1.5]
\setCJKmonofont{Noto Sans SC}
\setmonofont{Consolas}[Scale=0.9]

% ---------- 版式 ----------
\setlength{\headheight}{15pt}
\geometry{left=2.7cm, right=2.5cm, top=2.7cm, bottom=2.5cm,
          headsep=0.6cm, footskip=0.9cm}
\pagestyle{fancy}
\fancyhf{}
\renewcommand{\headrulewidth}{0.4pt}
\fancyhead[LE]{\nouppercase{\leftmark}}
\fancyhead[RO]{\nouppercase{\rightmark}}
\fancyfoot[C]{\thepage}
\fancypagestyle{plain}{\fancyhf{}\renewcommand{\headrulewidth}{0pt}\fancyfoot[C]{\thepage}}

\tolerance=1500
\emergencystretch=3em
\hfuzz=2pt

% ---------- 定理类环境（编号与原书一致：章.节.序号） ----------
\theoremstyle{plain}
\newtheorem{theorem}{定理}[section]
\newtheorem{lemma}[theorem]{引理}
\newtheorem{proposition}[theorem]{命题}
\newtheorem{corollary}[theorem]{推论}
\theoremstyle{definition}
\newtheorem{definition}[theorem]{定义}
\newtheorem{example}[theorem]{例}
\newtheorem{exercise}[theorem]{习题}
\theoremstyle{remark}
\newtheorem{remark}[theorem]{注}

\renewcommand{\thetheorem}{\thesection.\arabic{theorem}}
% 公式按章连续编号，如 (2.13)
\counterwithin*{equation}{chapter}
\renewcommand{\theequation}{\thechapter.\arabic{equation}}

% ---------- 校注环境 ----------
\newenvironment{annote}{%
  \par\vspace{0.4em}\begingroup\small
  \noindent\textcolor{black!45}{\rule[-0.25em]{2.2pt}{1.1em}}\hspace{0.5em}%
  \textsf{\textbf{译注}}\hspace{0.4em}%
}{\par\endgroup\vspace{0.4em}}

% ---------- 常用记号 ----------
\newcommand{\R}{\mathbb{R}}
\newcommand{\C}{\mathbb{C}}
\newcommand{\Z}{\mathbb{Z}}
\newcommand{\Q}{\mathbb{Q}}
\newcommand{\PP}{\mathbb{P}}
\newcommand{\CP}{\mathbb{CP}}
\newcommand{\RP}{\mathbb{RP}}
\newcommand{\calO}{\mathcal{O}}
\newcommand{\Fix}{\operatorname{Fix}}
\newcommand{\Aut}{\operatorname{Aut}}
\newcommand{\Inn}{\operatorname{Inn}}
\newcommand{\Out}{\operatorname{Out}}
\newcommand{\Ad}{\operatorname{Ad}}
\newcommand{\Hom}{\operatorname{Hom}}
\newcommand{\Id}{\operatorname{Id}}
\newcommand{\GL}{\operatorname{GL}}
\newcommand{\U}{\operatorname{U}}
\newcommand{\Spm}{\operatorname{Spm}}
\newcommand{\Cov}{\operatorname{Cov}}
\newcommand{\Fib}{\operatorname{Fib}}
\newcommand{\Gal}{\operatorname{Gal}}
\newcommand{\Cl}{\operatorname{Cl}}
\newcommand{\restr}[2]{#1|_{#2}}
\newcommand{\inputif}[1]{\IfFileExists{chapters/#1.tex}{\input{chapters/#1}}{}}

% ---------- 第 1 章专用 ----------
% 图版从略：原位留一行「（图 1.x 从略：图注译文）」
\newcommand{\figomit}[2]{\par\phantomsection\label{fig:#1}%
  \begin{center}\small（图 #1 从略：#2）\end{center}\par}
% PGL^(N)^♯（余扩张）与 PGL^#(N)（中心扩张）
\newcommand{\PGLs}{\operatorname{PGL}^{\sharp}(N)}
\newcommand{\PGLh}{\operatorname{PGL}^{\#}(N)}

\begin{document}

\frontmatter
\begin{titlepage}
\centering
\vspace*{2.6cm}
{\Huge\bfseries 模空间的几何与量子化\par}
\vspace{1.1cm}
{\Large Geometry and Quantization\\[0.2em] of Moduli Spaces\par}
\vspace{2.2cm}
{\large Vladimir Fock \quad Andrey Marshakov\\[0.35em]
Florent Schaffhauser \quad Constantin Teleman\\[0.35em]
Richard Wentworth\par}
\vspace{0.8cm}
{\large 著\par}
\vspace{1.2cm}
{\large 编者\\[0.3em]
Luis Álvarez-Cónsul \quad Jørgen Ellegaard Andersen \quad Ignasi Mundet i Riera\par}
\vspace{1.6cm}
{\large 中文全译本\par}
\vfill
{\small 底本：\emph{Geometry and Quantization of Moduli Spaces}, Advanced Courses in
Mathematics -- CRM Barcelona, Birkhäuser, 2016（ISBN 978-3-319-33577-3，\\
DOI 10.1007/978-3-319-33578-0）。全书 219 印刷页，四章；本卷四章全译，每章末附该章参考文献。\par}
\end{titlepage}

\chapter*{译者说明}
\addcontentsline{toc}{chapter}{译者说明}
\markboth{译者说明}{译者说明}

\section*{一、底本与范围}

本卷译自 \emph{Geometry and Quantization of Moduli Spaces}（Advanced Courses in
Mathematics -- CRM Barcelona，Birkhäuser，2016；ISBN 978-3-319-33577-3，
DOI 10.1007/978-3-319-33578-0）。原书汇集 2012 年 3 月至 6 月在巴塞罗那自治大学
数学研究中心（CRM，Bellaterra）举办的同名学期课程的四门讲义，共四章：

\begin{enumerate}[label=\arabic*., leftmargin=2.2em, itemsep=0.15em]
\item V.~Fock、A.~Marshakov，《Loop Groups, Clusters, Dimers and Integrable
      Systems》，印刷页 1--65；
\item F.~Schaffhauser，《Lectures on Klein Surfaces and Their Fundamental Group》，
      印刷页 67--107；
\item C.~Teleman，《Five Lectures on Topological Field Theory》，印刷页 109--163；
\item R.~A.~Wentworth，《Higgs Bundles and Local Systems on Riemann Surfaces》，
      印刷页 165--215。
\end{enumerate}

本卷是试点卷，收入\textbf{前言}与\textbf{第 2 章}的完整译文。第 1、3、4 章已建
骨架与任务书（见工程文件 \texttt{work/notes/}），译文待后续各卷补齐。全书导读见
下一节《全书导读》，其中第 2 章之外的概述均据原书前言与各章导言写成，仅作导航，
不含译文。

\section*{二、页码对照与引用}

原书正文印刷页 $=$ 底本 PDF 页 $-10$：第 2 章起于印刷页 67（PDF 页 77），止于
印刷页 107（PDF 页 117）；参考文献在印刷页 107--108。译文中的定理、公式编号
（如定义 2.1.1、式 (2.13)、习题 2.2.30）与原文完全一致，可直接回查底本；正文中的
「原书 p.~$n$」一律指印刷页码。

\section*{三、体例}

\begin{enumerate}[label=\arabic*., leftmargin=2.2em, itemsep=0.15em]
\item 定理类环境按原书体系编号为「章.节.序号」，公式按章连续编号，均不改动。
\item 术语在首次出现处给出原文，其后统一用中文；全书的对照见 \texttt{成品/术语对照表.md}。
\item 原文出现笔误、或原文自身前后不一致处，一律\textbf{照译}，另以「译注」说明，
      不静默改正。
\item 原文的习题全数保留并翻译；原书无插图者，此处亦不补图。
\item 底本 PDF 的文本层存在连字（\texttt{fl}、\texttt{fi}）与重音符号的私用编码
      （如把 $\tilde\sigma$ 抽成不可见控制符），凡公式一律以页面图像逐符号核对后再录入，
      不采信文本层。
\end{enumerate}

\section*{四、已知存疑}

\begin{itemize}[leftmargin=1.6em, itemsep=0.2em]
\item 原书前言称第 3 章作者 C.~Teleman 时为「University of Oxford」，而封面后的作者
      单位页写「University of California, Berkeley」。两处皆照录，见译注。
\item 版权页作者单位一栏的文本层行序错乱（Marshakov 的单位被排到 Schaffhauser 之后），
      录入时按封面与前言校正为：Andrey Marshakov，ITEP 与 Lebedev 物理研究所，莫斯科。
\end{itemize}

\begin{annote}
本卷为《数译》工程 SY05-30 的试点产物，目的是把「通读—术语—翻译—编译—资料卡—
插件反馈」整条管线走通并留下可复核的记录。工程文件（通读笔记、疑点清单、术语表、
编译说明、校订日志、插件测试笔记、工作流记录）随源码一并入库，供后续各章复用。
\end{annote}

\chapter*{前言}
\addcontentsline{toc}{chapter}{前言}
\markboth{前言}{前言}

本书基于「模空间的几何与量子化」（The Geometry and Quantization of Moduli Spaces）
学期课程期间开设的四门高级课程，该学期课程于 2012 年 3 月至 6 月在巴塞罗那
Bellaterra 的数学研究中心（Centre de Recerca Matemàtica，CRM）举行。模空间除了
在数学的许多领域中扮演重要角色之外，近几十年来还成为理解高能物理中若干现象的关键，
并因此促成了数学与物理之间的互动——这种互动对两门学科而言都可谓收获惊人。我们希望
读者能在本卷中欣赏并享受这种互动的美。

这些课程所处理的几个专题，在模空间这一广阔研究领域内正经历非同寻常的增长，它们也
反映了该研究学期中反复出现的一条线索，即黎曼面上局部系统（local system）的模空间，
无论是从经典的角度还是从量子的角度。

第 1 章由 V.~V.~Fock（法国斯特拉斯堡，IRMA）与 A.~Marshakov（俄罗斯莫斯科，ITEP）
撰写，题为《圈群、cluster、二聚体与可积系统》（Loop Groups, Clusters, Dimers and
Integrable Systems），属于 cluster 可积系统的范畴；更确切地说，属于 A.~B.~Goncharov
与 R.~Kenyon 借助二维环面上二部图的统计物理中二聚体模型（dimer model）理论所构造的
那一类可积系统。这篇研究论文描述了这些可积系统与在适当的仿射 Poisson--Lie 群上构造的
标准可积系统之间的一个同构，从而既给出了前者的新构造，也给出了后者的新视角。本章包含
非常有趣的例子，并指出了若干令人兴奋的推广方向。

第 2 章由 F.~Schaffhauser（哥伦比亚波哥大，Universidad de los Andes）撰写，题为
《克莱因曲面及其基本群讲义》（Lectures on Klein Surfaces and Their Fundamental
Group），是关于克莱因曲面的基本群的一篇综述性讲义。它包括克莱因曲面与实代数曲线的
入门、对克莱因曲面基本群的细致处理，以及对克莱因曲面基本群的线性表示与酉表示、
酉表示簇的讨论。许多定义与结果都以简单而有描述力的例子作动机或说明，并且每一节都配有
有趣的习题，读者可用以巩固所学。

接下来是第 3 章，由 C.~Teleman（英国牛津，牛津大学）撰写，题为《拓扑场论五讲》
（Five Lectures on Topological Field Theory）。它从 M.~F.~Atiyah 整理成体系的
拓扑量子场论（TQFT）经典定义出发，经由不同的例子，引出这些理论如今在多数情形下已知
具有的额外结构——正如 J.~Lurie 最近证明的所谓配边猜想（cobordism hypothesis）所
预言的那样。作者并未给出扩展 TQFT 的完全详细的定义（读者可在 Lurie 的论文中找到），
而是借助精心挑选的例子，阐释这一观念实质的各个侧面，以及证明其相容性与配边猜想所面临
的困难。在我们看来，这使本文成为 Lurie 论文的一份极佳 companion。
\begin{annote}
「companion」在原文中即英文词，指与之配套、可对读的著作。此处照录，未译作「伴侣」。
\end{annote}

第四也是最后一章，由 R.~A.~Wentworth（美国 College Park，马里兰大学）撰写，题为
《黎曼面上的 Higgs 丛与局部系统》（Higgs Bundles and Local Systems on Riemann
Surfaces）。Higgs 丛由 N.~Hitchin 在三十年前引入，现已在不同领域中证明具有中心地位，
例如低维拓扑、代数几何、数学物理，乃至数论（比如几年前 B.~C.~Ngô 对 Langlands 纲领中
基本引理的证明）。本章对这一概念作了教学式的入门介绍，但限于黎曼面的情形（这正是
Hitchin 论文中的原始语境）。它给出了 Higgs 丛模空间与局部系统模空间之间对应的一个
近乎自足的证明，然后将其应用于研究 oper 模空间——这一对象在几何 Langlands 纲领近来的
发展中起着关键作用。

我们要感谢本书各位作者的重大贡献，既为准备了一套非常有趣的讲义，也为撰写了本册笔记。
我们也感谢各位审稿人，他们毫不妥协地帮助了我们。

「模空间的几何与量子化」学期课程得以举办，有赖于不同机构的慷慨资助。为此我们要感谢
数学研究中心（CRM）、通过研究网络 ITGP（Interactions between Topology, Geometry
and Physics，拓扑、几何与物理之间的相互作用）给予支持的欧洲研究基金会、IMUB
（Institut de Matemàtiques de la Universitat de Barcelona，巴塞罗那大学数学研究所）
以及 SCM（Societat Catalana de Matemàtiques，加泰罗尼亚数学学会）。最后，我们感谢
CRM 的工作人员，他们在课程筹备与运行期间的效率与帮助令人感念。

\vspace{1.2em}
\begin{flushright}
马德里，Luis Álvarez-Cónsul（ICMAT）\\
奥胡斯，Jørgen Ellegaard Andersen（QGM）\\
巴塞罗那，Ignasi Mundet i Riera（UB）
\end{flushright}

\chapter*{全书导读}
\addcontentsline{toc}{chapter}{全书导读}
\markboth{全书导读}{全书导读}

\noindent\textbf{一句话主旨。}全书讲的是黎曼面上局部系统（local system）模空间的两副面孔：
一副是「经典」的——基本群的表示簇、Higgs 丛与 Hitchin 系统；一副是「量子」的——圈群、
可积系统、拓扑场论与几何量子化。四章各取一面，而第 2 章给出的 $\pi_1^{\R}$ 及其表示簇，
正是另三章反复用到的地基。

\section*{第 1 章\quad 圈群、cluster、二聚体与可积系统}

作者按「两种构造、一个同构」来组织：先回顾仿射 Poisson--Lie 群（$\hat A_{N-1}$ 型，
即 $PGL(N)$ 的圈群）上的标准可积系统，再回顾 Goncharov--Kenyon 用二维环面上的二部
「二聚体」（dimer）模型造出的 cluster 可积系统，然后证明两者同构。中段的
\S1.3--\S1.4 用 double Bruhat 胞腔的 cluster 参数化把简单群与圈群统一处理，关键工具是
Cartan--Weyl 生成元、coextended 仿射 Weyl 群与 Thurston 图；\S1.5 把 dimer 分割函数
与离散 Dirac 算子联系起来，得到谱子流形与 Poisson 括号；\S1.8 的例子里有最简单相对论
Toda 链、rank two 相对论 Toda 链，以及五角星映射（pentagram map）。六个附录补足了
cluster 簇、Thurston 图与 Schwartz 坐标的技术细节。

\emph{读法提示：}本章公式密集、图多，宜从 \S1.1 的三小节读起，先抓住「同一个可积系统
的两种坐标」这一主线，再进 \S1.5 的 dimer 机器；附录不必顺读，用到再查。

\section*{第 2 章\quad 克莱因曲面及其基本群（本卷全译）}

三节：\S2.1 讲 Klein 曲面的一般理论（实代数曲线、拓扑分类、双解析结构），\S2.2 定义
克莱因曲面的基本群 $\pi_1^{\R}((X,\sigma);x)$ 并证明它分类实覆盖，\S2.3 讨论线性表示、
酉表示与酉表示簇的实结构。核心是一条短正合列
\[
1\longrightarrow \pi_1(X;x)\longrightarrow \pi_1^{\R}((X,\sigma);x)
\longrightarrow \operatorname{Gal}(\C/\R)\longrightarrow 1,
\]
以及一个类比：SGA-1 意义下代数基本群之于拓扑基本群，正如 $\pi_1^{\R}$ 之于 $\pi_1$。
每一节都配了习题（2.1.17--2.1.27、2.2.22--2.2.31、2.3.5--2.3.10），其中不少是正文结论的
另一证明或具体算例，值得动手。

\emph{读法提示：}若只关心「实结构如何作用在表示簇上」，可直接读 \S2.2.3 与 \S2.3.2，
再回补 \S2.2.2 的构造；定理 2.1.8（Weichold）是全章唯一的分类定理，\S2.1.2 的
$(g,k,a)$ 三元组要记牢。

\section*{第 3 章\quad 拓扑场论五讲}

从 Atiyah 的 TQFT 公理出发（\S3.1），经有限群规范理论、Frobenius 代数、二维
Yang--Mills、上同调场论等例子，过渡到扩展 TQFT 与高阶范畴（\S3.3），再在一维、二维
情形陈述并「看图」证明配边猜想（\S3.4，含对偶对象、Serre 扭、有向与 $r$-spin 理论、
有向把柄的伴随），最后走向高维（\S3.5，含 Drinfeld 中心、融合范畴与 Serre 自同构）。
全章的主线是：一个 $n$ 维扩展 TQFT 的「值」由一点处的对象加上可对偶性决定。

\emph{读法提示：}这是 Lurie 论文的 companion，读法也应如是——先把 \S3.4 的二维情形
吃透（Frobenius 代数与配边猜想的对应最干净），再看 \S3.5 如何把同一套语言升到高维。

\section*{第 4 章\quad 黎曼面上的 Higgs 丛与局部系统}

\S4.2 从 Dolbeault 面讲 Higgs 丛与模空间，并给出 Hitchin--Kobayashi 对应（稳定
Higgs 丛 $\leftrightarrow$ 规范场方程解）；\S4.3 从 Betti 面讲表示簇、局部系统与全纯
联络，给出 Corlette--Donaldson 定理与 hyperkähler 约化，于是得到两个模空间之间的
对应（非阿贝尔 Hodge 对应）；\S4.4 把对应用到微分方程：单值化、高阶方程、opers 与
Eichler--Shimura 同构，收束于几何 Langlands 的入口。

\emph{读法提示：}若已熟悉 Hitchin 系统，可从 \S4.3.3 切入；若想看清「对应」的解析
机制，则 \S4.2.3 与 \S4.3.4 是两根支柱，值得逐行核。

\section*{四章之间的接口}

\begin{itemize}[leftmargin=1.8em, itemsep=0.25em]
\item 第 2 章的 $\pi_1^{\R}$ 与酉表示簇的实结构，是第 4 章「实 Higgs 丛 / 实局部系统」
      的语言基础；第 2 章定理 2.3.4（$\Phi$ 的像落在 $\Fix(\kappa)$）正是「实点」这一
      观念的拓扑原型。
\item 第 1 章的 dimer/cluster 机器与第 3 章的二维 TQFT 共享「二维环面 + 二部图 +
      分割函数」这套组合学，\S1.5.6 的谱子流形与 \S3.3 的融合范畴值得对照。
\item 第 3 章的扩展 TQFT 与第 4 章的几何 Langlands 在「oper 模空间 = 某些 TQFT 的
      状态空间」这一点上相遇，这是当前研究的前沿接口之一。
\end{itemize}

\begin{annote}
本节只作导航：第 1、3、4 章的内容概述据原书前言与目录写成，本卷未收其译文，
引用时请以原书页码为准（第 1 章 p.1--65，第 3 章 p.109--163，第 4 章 p.165--215）。
\end{annote}

\tableofcontents

\mainmatter
% 章号与原书一致：第 1 章 Fock–Marshakov，第 2 章 Schaffhauser。
\setcounter{chapter}{0}
\inputif{ch01a}
\inputif{ch01b}
\inputif{ch01c}
\inputif{ch01d}
\inputif{ch01e}
\inputif{ch01f}
\inputif{ch01g}
\inputif{bib01}

\setcounter{chapter}{1}
\inputif{ch02a}
\inputif{ch02b}
\inputif{ch02c}
\inputif{ch02d}
\inputif{ch02e}
\inputif{ch02f}
\inputif{ch02g}
\inputif{ch02h}
\inputif{ch02i}
\inputif{ch02j}

\backmatter
\renewcommand{\bibname}{第 2 章参考文献}

\begin{thebibliography}{99}\setlength{\itemsep}{0.25em}
\bibitem{alling} N.~L.~Alling and N.~Greenleaf, \emph{Foundations of the Theory of Klein
Surfaces}, Lecture Notes in Mathematics 219, Springer-Verlag, Berlin, 1971.
\bibitem{atiyah} M.~F.~Atiyah, K-theory and reality, \emph{Quart. J. Math. Oxford}
\textbf{17} (1966), 367--386.
\bibitem{conway} J.~B.~Conway, \emph{Functions of One Complex Variable}, Graduate Texts in
Mathematics 11, Springer-Verlag, New York (second edition), 1978.
\bibitem{donaldson} S.~Donaldson, \emph{Riemann Surfaces}, Oxford Graduate Texts in
Mathematics 22, Oxford University Press, Oxford, 2011.
\bibitem{forster} O.~Forster, \emph{Lectures on Riemann Surfaces}, Graduate Texts in
Mathematics 81, Springer-Verlag, New York, 1991; translated from the 1977 German original
by B.~Gilligan, reprint of the 1981 English translation.
\bibitem{grothendieck} A.~Grothendieck, \emph{Revêtements étales et groupe fondamental
(SGA-1)}, Documents Mathématiques 3, Société Mathématique de France, Paris, 2003;
Séminaire de Géométrie Algébrique du Bois Marie 1960--61, directed by A.~Grothendieck,
with two papers by M.~Raynaud, updated and annotated reprint of the 1971 original Lecture
Notes in Math.~224, Springer, Berlin.
\bibitem{harnack} A.~Harnack, Ueber die vieltheiligkeit der ebenen algebraischen curven,
\emph{Math. Ann.} \textbf{10}(2) (1876), 189--198.
\bibitem{huisman1} J.~Huisman, Real Teichmüller spaces and moduli of real algebraic
curves, in \emph{Real Algebraic Geometry and Ordered Structures}, Contemp. Math.~253
(1996), 145--177; Amer. Math. Soc., Providence, RI, 2000.
\bibitem{huisman2} J.~Huisman, The equivariant fundamental group, uniformization of real
algebraic curves, and global complex analytic coordinates on Teichmüller spaces,
\emph{Ann. Fac. Sci. Toulouse Math.} \textbf{10}(4) (2001), 659--682.
\bibitem{liu} C.~C.~M.~Liu, \emph{Moduli of J-holomorphic curves with Lagrangian boundary
conditions}, ProQuest LLC, Ann Arbor, MI, 2002; thesis, Harvard University.
\bibitem{massey} W.~S.~Massey, \emph{A Basic Course in Algebraic Topology}, Graduate Texts
in Mathematics 127, Springer-Verlag, New York, 1991.
\bibitem{milnor} J.~Milnor, On the Betti numbers of real varieties, \emph{Proc. Amer.
Math. Soc.} \textbf{15} (1964), 275--280.
\bibitem{mochizuki} S.~Mochizuki, Topics surrounding the anabelian geometry of hyperbolic
curves, in \emph{Galois Groups and Fundamental Groups}, Math. Sci. Res. Inst. Publ.~41,
119--165; Cambridge Univ. Press, Cambridge, 2003.
\bibitem{natanzon} S.~M.~Natanzon, Klein surfaces, \emph{Uspekhi Mat. Nauk}
\textbf{45}(6) (1990), 47--90; English version: \emph{Russian Mathematical Surveys}
\textbf{45}(5) (1990), 53--108.
\bibitem{ratcliffe} J.~G.~Ratcliffe, \emph{Foundations of Hyperbolic Manifolds}, Graduate
Texts in Mathematics 149, Springer, New York (second edition), 2006.
\bibitem{schaffhauser} F.~Schaffhauser, Real points of coarse moduli schemes of vector
bundles on a real algebraic curve, \emph{J. Symplectic Geom.} \textbf{10}(4) (2012),
503--534.
\bibitem{scott} P.~Scott, The geometries of 3-manifolds, \emph{Bull. London Math. Soc.}
\textbf{15}(5) (1983), 401--487.
\bibitem{szamuely} T.~Szamuely, \emph{Galois Groups and Fundamental Groups}, Cambridge
Studies in Advanced Mathematics 117, Cambridge University Press, Cambridge, 2009.
\bibitem{thom} R.~Thom, Sur l'homologie des variétés algébriques réelles, in
\emph{Differential and Combinatorial Topology} (a symposium in honor of Marston Morse),
255--265; Princeton Univ. Press, Princeton, N.J., 1965.
\bibitem{weichold} G.~Weichold, \emph{Über symmetrische Riemannsche Flächen und die
Periodizitätsmodulen der zugerhörigen Abelschen Normalintegrale erster Gattung}, 1883,
Leipziger Dissertation.
\end{thebibliography}


\end{document}
